\documentclass[10pt,a4paper]{amsart}
\usepackage[margin=19mm]{geometry}
\usepackage{amsmath,amssymb,mathtools}
\usepackage[scr=rsfs,cal=euler]{mathalfa}
\usepackage{xfrac}
\usepackage[unicode,hidelinks,bookmarks=false]{hyperref}
\newcommand{\ie}{i.e.\ }
\newcommand{\cf}{cf.\ }
\newcommand{\resp}{respectively\ }
\newcommand{\Subsection}{\subsection}
\providecommand{\nobreakdash}{\nobreak\hbox{-}}
\setlength{\emergencystretch}{2em}
\multlinegap0pt
\usepackage{bm}
\usepackage{bbm}
\usepackage{dsfont}
\usepackage{xspace}
\usepackage{cancel}
\usepackage{mathtools}
\usepackage{amsthm}
\makeatletter
\@ifundefined{theo}{\newtheorem{theo}{Theorem}}{}
\@ifundefined{prop}{\newtheorem{prop}[theo]{Proposition}}{}
\makeatother
\newcommand\da{\underline{\mbox{\rm dim}}_{\,\theta}} % lower theta dimension
\newcommand\lbd{\underline{\mbox{\rm dim}}_{\rm B}\,} % lower box dimension
\newcommand\ubd{\overline{\mbox{\rm dim}}_{\rm B}\,} % upper box dimension
\newcommand\uda{\overline{\mbox{\rm dim}}_{\,\theta}} % upper theta dimension
\renewcommand{\epsilon}{\varepsilon}
\title[Intermediate dimensions of complementary sets]{Intermediate dimensions of complementary sets}
\author{Nicolas Angelini, Ursula Molter}
\date{Source: arXiv:2412.12999v3 (CC BY 4.0)}
\begin{document}
\maketitle

\noindent\emph{Source and licence.} Intermediate dimensions of complementary sets. Version arXiv:2412.12999v3 (CC BY 4.0), distributed under the Creative Commons Attribution 4.0 International licence, as recorded in the arXiv licence metadata for this exact version and shown on the abstract page https://arxiv.org/abs/2412.12999v3. Full rights notice: licenses/arxiv-2412.12999-CC-BY-4.0.txt.\par\smallskip

\noindent\textit{Editorial scope.} Counted unit: the Theo whose source label is \texttt{Maintheo} in the article (source0.tex lines 521--535, source label \texttt{Maintheo}), delivered together with the article's own complete original proof of it (source0.tex lines 537--779, one \texttt{proof} environment of 243 lines). No mathematics is written, repaired or re-derived: statement and proof are the article's own text line for line. Required context retained uncounted: hausdorff cantor (lines 196--202); dimension de suceciones (lines 382--391). Every definition, display and label of those lines is retained unchanged. Omitted from this package and not redistributed: 1--195, 203--381, 392--520, 536--536, 780--935. Nothing inside a retained excerpt is omitted or altered: every retained body is delivered exactly as the article's lines are written, so no omission receipt of any kind accompanies this package. Citations: the retained excerpts carry no citation command at all, so no bibliography entry of the article is transcribed and no reference is redistributed. Reference adaptation: none was needed and none was made; every reference inside the delivered unit resolves inside it, so no unresolved reference or citation is produced. Printed slips kept literal: no printed word, symbol, hypothesis or number of the counted statement or of its proof was altered. Layout and class: the article's own class is \texttt{amsart} with packages including amsmath, amsthm, amscd, amsfonts, amssymb, graphicx, color; the wrapper is the lane's pinned \texttt{amsart} editorial wrapper at 10pt on A4, which restates the theorem-like environments the retained text needs and adds the article's own macro definitions that the retained text needs (\texttt{\textbackslash da}, \texttt{\textbackslash epsilon}, \texttt{\textbackslash lbd}, \texttt{\textbackslash ubd}, \texttt{\textbackslash uda}), with extra wrapper packages (bm, bbm, dsfont, xspace, cancel, mathtools). The delivered numbering is the unit's own. Assets: no retained span contains a figure environment, a graphics-inclusion command, a PDF-page inclusion, a table or a TikZ picture; no image file is copied into this package, the package declares no asset dependency and no image file is redistributed with it. Licence: version arXiv:2412.12999v3 of this article is distributed under the Creative Commons Attribution 4.0 International licence, as recorded in the arXiv licence metadata for this exact version and shown on the abstract page https://arxiv.org/abs/2412.12999v3. A full rights notice is delivered at licenses/arxiv-2412.12999-CC-BY-4.0.txt. Characters: every retained excerpt is pure ASCII, so the delivered engine, XeLaTeX with its default Latin Modern text font, reports no missing character.
\begin{prop}\label{hausdorff cantor}
Let $a = \{a_n\}$ be a positive summable sequence, and let $C_a$ be the decreasing Cantor set associated with $a$. Then
\[
\dim_H C_{a} = \liminf_{n \to \infty} \frac{\log n}{-\log (x_n/n)}.
\]
Furthermore, any set $E \in \mathcal{C}_{a}$ satisfies $\dim_H E \leq \dim_H C_{a}$. Conversely, given any $s \in [0, \dim_H C_{a}]$, there exists a set $E \in \mathcal{C}_a$ such that $\dim_H E = s$.
\end{prop}

\begin{theo}\label{dimension de suceciones} 
Let \(a=\{a_n\}_{n\in\mathbb{N}}\) be a non increasing summable sequence. Then for all \(\theta\in [0,1]\), we have
\begin{equation}\label{Theo inciso 11}
\uda D_{\textbf{a}} = \frac{\theta \, \ubd D_{\textbf{a}} }{1 - (1-\theta)\, \ubd D_{\textbf{a}}}
\end{equation}
and  
\[
\da D_{\textbf{a}} = \frac{\theta \, \lbd D_{\textbf{a}} }{1 - (1-\theta)\, \lbd D_{\textbf{a}} }.
\]
\end{theo}

\begin{theo}\label{Maintheo}
Let \( \theta \in (0,1] \) and let \( a = \{a_n\} \) be a positive, non-increasing, summable sequence. If 
\[
2 < \inf_n \frac{s_n(a)}{s_{n+1}(a)} \leq \sup_n \frac{s_n(a)}{s_{n+1}(a)} < \infty,
\]
then for every 
\[
t \in \left[ \overline{\dim}_\theta D_a, \, \overline{\dim}_\theta C_a \right],
\]
there exists a set \( E = E(t, \theta) \in \mathcal{C}_a \) such that \( \overline{\dim}_\theta E = t \). An analogous result holds for the upper $\theta$-intermediate dimension, i.e., for every 
\[
t \in \left[ \underline{\dim}_\theta D_a, \, \underline{\dim}_\theta C_a \right],
\]
there exists a set \( E' = E'(t, \theta) \in \mathcal{C}_a \) such that \( \underline{\dim}_\theta E' = t \).
\end{theo}

\begin{proof}
The proof proceeds by explicit construction of the sets $E$ and $E'$. Given the sequence \(a\), we will construct a Cantor-type set \(C_b\) associated to a subsequence \(b\) of \(\{a_n\}\), and a countable set \(D_{a'}\) associated to the remaining elements of \(\{a_n\}\). We will then define
\[
E = C_b \cup D_{a'},
\]
and show that \( \overline{\dim}_\theta E = t\).

First, we note that the hypothesis $2<\inf_n\frac{s_n(a)}{s_{n+1}(a)}\leq \sup_n\frac{s_n(a)}{s_{n+1}(a)}<\infty $ implies that $a$ is a doubling sequence.

Let \( \theta \in (0,1] \) and fix \( t \in ( \overline{\dim}_\theta D_a , \overline{\dim}_\theta C_a) \). Define
\[
s := \frac{\overline{\dim}_\theta C_a}{t} > 1.
\]
For each \(n\), let \(j=j(n)\) be such that
\begin{equation}\label{j(n)}
    a_{2^{j+1}} < (a_{2^n})^s \leq a_{2^j}, \qquad j \geq n.
\end{equation}


\[ b_{2^k+i} = \begin{cases}                  a_{2^{j(k)+1}+i} & \text{if $j(k)>j(k-1)$,}\\                  a_{2^{j(k)+1}+\sum_{r=1}^m 2^{k-r}+i} &    \text{if $j(k)=\cdot\cdot\cdot=j(k-m)>j(k-m-1)$,}           \end{cases} \]

for \(k\ge 1\) and \(0\le i<2^k\), where \(m\) denotes the largest integer with \(0\le m\le k-1\) such that \(j(k)=j(k-1)=\cdots=j(k-m)\). 
We also adopt the convention \(j(0)=0\).\\
By construction and using the doubling condition of $a$, there exists a constant $c>0$ such that for each \(k \geq 1\) and \(0 \leq i < 2^k\),
\[
c^{-1} (a_{2^k})^s \leq b_{2^k+i} \leq c (a_{2^k})^s.
\]


Now, observe that

$$
\inf_n \frac{s_n(b)}{s_{n+1}(b)} > 2.
$$

Indeed, recall that

$$
s_n(a) = 2^{-n} \sum_{i=2^n}^{2^{n+1}-1} a_i + 2 s_{n+1}(a).
$$

From this identity, the hypothesis

$$
2 < \inf_n \frac{s_n(a)}{s_{n+1}(a)}
$$

can be rewritten in the equivalent form

$$
\sup_n \frac{s_{n+1}(a)}{r_n(a)} < \infty,
$$

where

$$
r_n(a) = 2^{-n}\sum_{i=2^n}^{2^{n+1}-1} a_i.
$$

Using the doubling condition for $a$, one directly obtains

$$
\frac{s_{n+1}(b)}{r_n(b)}\leq \frac{C}{2^{-(n+1)}} \sum_{i=2^{n+1}}^{\infty}\left(\frac{a_i}{a_{2^n}}\right)^s
< C\, \frac{2^{-(n+1)}\sum_{i=2^{n+1}}^{\infty}a_i}{a_{2^n}}
\leq C'\frac{s_{n+1}(a)}{r_n(a)},
$$

for some constant $C$ and $C'$, and hence $\sup_n \tfrac{s_{n+1}(b)}{r_n(b)} < \infty$.


It follows that for doubling decreasing sequences with this property we have\[b_{2^n}\asymp s_n(b).\]Indeed, if \(s_{n+1}(b)/s_n(b)<1/2-\delta\), then\[2\delta s_n(b) \leq s_n(b)-2s_{n+1}(b)=2^{-n}\sum_{j=2^n}^{2^{n+1}-1}b_j\leq b_{2^{n}}\leq c\, 2^{-n} \sum_{j=2^n}^{2^{n+1}-1}b_j \leq c\,s_n(b).\]Hence\[s_n(b)\asymp b_{2^n}\asymp (a_{2^n})^s\asymp (s_n(a))^s.\]


Now, we will prove that $\overline{\dim}_\theta C_b=s^{-1}\overline{\dim}_\theta C_a.$

Let $d > \overline{\dim}_\theta C_a$. Then, for all $\epsilon > 0$, there exists $\delta_0 > 0$ such that for every $\delta \in (0,\delta_0)$ there is a cover $\{U_i\}$ of $C_a$ satisfying
\begin{equation}\label{ecuacion definicion}
\delta \leq |U_i| \leq \delta^\theta
\quad \text{and} \quad
\sum_i |U_i|^d < \epsilon.
\end{equation}
Let $\delta$ and the cover $\{U_i\}$ satisfy \eqref{ecuacion definicion}. For each $i$, choose $k(i)$ so that
\[
s_{k(i)+1}(a) \leq |U_i| < s_{k(i)}(a).
\]

It follows that each $U_i$ can intersect at most four closed intervals from step $k(i)-2$ in the construction of $C_a$, namely the intervals $I_{j_n}^{\,k(i)-2}(a)$ with $n=1,\dots,4$. Consequently, we can cover $C_a$ by
\[
C_a \subset \bigcup_i \bigcup_{n=1}^4 I_{j_n}^{\,k(i)-2}(a).
\]
Analogously, we can cover $C_b$ by selecting the corresponding intervals from the same steps in the construction of $C_b$, that is,
\[
C_b \subset \bigcup_i \bigcup_{n=1}^4 I_{j_n}^{\,k(i)-2}(b).
\]



Now, since
\[
\delta^s \leq c_1 s_{k(i)-3}(b) \leq c_1 |I_{j_n}^{\,k(i)-2}(b)| \leq c_1 s_{k(i)-1}(b) 
\leq c_1 c_2 \,(s_{k(i)-1}(a))^s \leq c_1 c_2\, |U_i|^s,
\]
for some constants $c_1, c_2$, it follows that
\[
|I_{j_n}^{\,k(i)-2}(b)| \in \bigl(c_1^{-1}\delta^s,\, c_1 c_2 \delta^{\theta s}\bigr).
\]
Moreover,
\[
\sum_i \sum_{n=1}^4 |I_{j_n}^{\,k(i)-2}(b)|^{d/s} < 4\, c_1 c_2 \,\epsilon,
\]
and therefore
\[
\overline{\dim}_\theta C_b \leq s^{-1}\,\overline{\dim}_\theta C_a.
\]
An analogous argument shows that
\[
s^{-1}\,\overline{\dim}_\theta C_a \leq \overline{\dim}_\theta C_b.
\]







We now define the complementary countable part. Let $\{\tilde{a}_n\}$ be a non-increasing sequence such that
\[
\{a_n : n \in \mathbb{N}\} = \{b_n : n \in \mathbb{N}\} \cup \{\tilde{a}_n : n \in \mathbb{N}\}.
\]
Define the associated countable set
\[
D_{a'} = \bigcup_{k=1}^{\infty} y_k, \qquad y_k = r + \sum_{i=k}^\infty \tilde{a}_i,
\]
where
\[
r = \sum_{i=1}^{\infty} b_i.
\]

Finally, set
\[
E = C_b \cup D_{a'}.
\]
Then \(E \in \mathcal{C}_a\). By finite stability and the invariance of the upper box dimension for complementary sets we have
\[
\ubd D_a =\overline{\dim}_B E = \max\{\overline{\dim}_B C_b, \overline{\dim}_B D_{a'}\}.
\]
Since
\[
\overline{\dim}_B C_b = s^{-1}\overline{\dim}_B C_a < \overline{\dim}_B D_a,
\]
we deduce that
\[
\overline{\dim}_B D_a = \overline{\dim}_B D_{a'}.
\]
Therefore
\[
\overline{\dim}_\theta D_a = \overline{\dim}_\theta D_{a'},
\]
by Theorem \ref{dimension de suceciones}.

Combining these results, we conclude that
\begin{align*}
\overline{\dim}_\theta E
&= \max\{s^{-1}\overline{\dim}_\theta C_a , \overline{\dim}_\theta D_{a'}\} \\
&= \max\{t, \overline{\dim}_\theta D_a\} = t,
\end{align*}
as required.


For the lower $\theta$-intermediate dimension, we must change strategy, as the lower $\theta$-intermediate dimension (as well as the lower box dimension) is not finitely stable. We will consider a similar set and proceed with a covering strategy to prove the theorem.

Let \( t \in (\underline{\dim}_\theta D_a, \underline{\dim}_\theta C_a) \). Proceed as before to construct a set \( E' \in \mathcal{C}_a \) such that \( E' = C_b \cup D_{a'} \), where $C_b$ is a decreasing Cantor set and \( D_{a'} \) is a countable set.

Since both sequences \( \{b_n\} \) and \( \{a_n\} \) are non-increasing, it follows from Proposition \ref{hausdorff cantor} and the subsequent remarks that \[ s^{-1}\underline{\dim}_\theta C_a = s^{-1}\underline{\dim}_B C_a = \underline{\dim}_B C_b = \underline{\dim}_\theta C_b  \] for all \( \theta \).

Now since  
\[
\frac{\theta \underline{\dim}_B C_a}{1 - \underline{\dim}_B C_a (1 - \theta)} = \underline{\dim}_\theta D_a < t = \frac{\underline{\dim}_B C_a}{s},
\]  
we have that  
\[
s < \frac{1 - (1 - \theta) \underline{\dim}_B C_a}{\theta}.
\]  
Therefore, we can choose \( t' > t \) and \( B > \underline{\dim}_B C_a > b \) such that  
\begin{equation}\label{condicion s}
t' > \max\left\{ \frac{B}{s},\ \frac{\theta B}{1 - B(1 - \theta)},\ \frac{\theta b - \theta(1 - b / t's)}{1 - b(1 - \theta) - (1 - b / t's)(1 - \theta)} \right\}
\end{equation}  
and  
\begin{equation}\label{condicion s 2}
s < \frac{1 - (1 - \theta) s t'}{\theta}.
\end{equation}

Now, since by construction of $E'$, $(s_n(a))^s \asymp s_n(b)$, it follows that there exist a subsequence $\{n_k\}$ and a constant $C > 0$ such that
\begin{equation}\label{Desigualdad de sna}
2^{-n_k(1/b)} \leq s_{n_k}(a) \leq 2^{-n_k(1/B)}
\end{equation}
and
\begin{equation}\label{Desigualdad de snb}
s_{n_k}(b) \leq C \, 2^{-n_k(s/B)}.
\end{equation}



Let $\delta_k := (s_{n_k}(a))^{1/(t' + \theta(1 - t'))}$. 
By equation \eqref{condicion s 2}, we have 
\[
(s_{n_k}(a))^s = \delta_k^{s t'(1 - \theta) + \theta s} \in (\delta_k, \delta_k^\theta),
\]
and therefore $C' s_{n_k}(b)$ is an admissible diameter for covering sets, for some constant $C'$.

Since the closed intervals of a decreasing Cantor set, say \( I_i^n \), satisfy \( s_{n+1} \leq |I_i^n| \leq s_{n-1} \), we can cover \( C_b \) with \( 2^{n_k + 1} \) intervals of length \(C' s_{n_k}(b) \).

To cover \( D_{a'} = \{x'_1, x'_2, \ldots\} \), we will cover \( \{x'_i : i = 1, \ldots, 2^{n_k}\} \) with \( 2^{n_k} \) sets of diameter \( \delta \), and the remaining points will be covered with
\[
\left\lceil \frac{2^{n_k} s_{n_k}(a'')}{\delta^\theta} \right\rceil \leq \left\lceil \frac{2^{n_k} s_{n_k}(a)}{\delta^\theta} \right\rceil \leq \frac{2^{-n_k(1 - t's)/t's}}{\delta^\theta} + 1,
\]
sets of diameter \( \delta^\theta \), using \eqref{Desigualdad de sna}.


Therefore we have created a cover \( \{U_i\} \), of \( E \), by sets whose diameters belong to \( [\delta, \delta^{\theta}] \), and that satisfies
\begin{align*}
    \sum_i |U_i|^{t'} &\leq 2C' 2^{n_k} s_{n_k}(b)^{t'} + 2^{n_k} \delta^{t'} + \left( c' \frac{2^{-n_k(1 - t's)/t's}}{\delta^\theta} + 1 \right) \delta^{\theta t'} \\
    &= S_1 + S_2 + S_3.
\end{align*}
Using \eqref{Desigualdad de snb} and \eqref{condicion s}, we have that \( S_1 \to 0 \) as \( n_k \to \infty \).

For \( S_2 \) we have
\begin{align*}
    2^{n_k} \delta^{t'} = 2^{n_k} (s_{n_k}(a))^{t'/(\theta + t'(1 - \theta))} \leq 2^{-n_k(t'/(B(\theta + t'(1 - \theta))) - 1)} \to 0
\end{align*}
since \( t' > \theta B / (1 - B(1 - \theta)) \).

On the other hand, \( S_3 \) will tend to \( 0 \) since by \eqref{condicion s} we have 
\[
t' > \frac{\theta b - \theta(1 - b / t's)}{1 - b(1 - \theta) - (1 - b / t's)(1 - \theta)}.
\]

Then, it follows that
\[
\underline{\dim}_\theta E' \leq t'
\]
and the result follows by letting \( t' \to t \) and noting that \( t \leq \underline{\dim}_\theta E' \) since \( C_b \subset E' \).
\end{proof}

\end{document}
